\section{What the repository statement adds}
The scope note for Family~179 documents a formalized classification of positive circulant sign Hadamard orders as exactly $1$ and $4$. It separately documents the even-length Barker consequence as lengths $2$ or $4$, and states that the paper's odd-length classification is outside that selected additional formal statement~\cite{oai179}. These are claims and scope descriptions from the repository, not results established by our finite experiment or by Proposition~\ref{thm:circulantnecessary}.

At this point you can understand the objects and the exact strength of the circulant statement. You also know a complete elementary supporting argument and its limit: excluding $4u^2$ for $u>1$ requires additional mathematics. Chapter~17 shows how to request that additional background without pretending that a short elementary argument has already supplied it.

\begin{lab}{Preventing an experiment from becoming a false proof}
\noindent\textbf{Deliberately flawed conclusion.} ``We found circulant examples only at lengths $1$ and $4$, so those are the only possible lengths.''

\noindent\textbf{Learner.} ``Separate the finite search statement, the necessary-order theorem, and the external classification. Which of these rules out order $16$? Which has actually been proved here?''

\noindent\textbf{Model audit.} The finite search stops at $13$, so says nothing about $16$. The necessary condition permits $16=4\cdot2^2$. The external classification excludes it, but its proof must be read or independently checked. The course has established neither that exclusion by enumeration nor the full classification by its elementary argument.
\end{lab}
